\begin{abstract}

Recent research has demonstrated the great potential of applying static analysis to hardware description languages (HDLs) for lightweight hardware bug detection, security analysis, and design understanding.
We observe that one fundamental aspect distinguishing hardware design from software programming is the explicit notion of temporal behavior with respect to clock cycles.
Developers' misunderstanding or incorrect implementation of intended temporal behavior can lead to temporal hardware bugs, such as signal asynchrony bugs, and security vulnerabilities, such as timing-flow information leaks.
However, existing research on hardware static analysis has not explored how to detect such temporal bugs and vulnerabilities or aid the understanding of hardware temporal behavior.

To address this gap, we present \emph{temporal influence analysis} (\tia), the first static analysis to reason about the temporal behavior of hardware designs written in Chisel---the most popular modern HDL.
\tia answers a fundamental question about hardware temporal behavior: after which clock cycles can each input to a circuit \emph{influence} each circuit location?
We formally define \tia and prove its soundness and convergence, as well as properties that help understand its precision.
To demonstrate \tia's potential as a foundation for diverse hardware analysis clients concerning temporal behavior, we develop three client analyses on top of it: detecting signal asynchrony bugs, identifying timing-flow information leaks, and understanding timeline interfaces.


Our evaluation of \tia focuses on two aspects:
(1) For \tia itself, we demonstrate that it efficiently analyzes 9 popular real-world Chisel projects, averaging 2.2K GitHub stars and totaling 5.9M LoC, within 50 minutes in total.
\tia achieves high recall (95\% on average) and fine precision (60\% on average) in capturing temporal influence information in these projects.
(2) For the three client analyses built on \tia, we demonstrate that they successfully capture all signal asynchrony bugs, timing-flow vulnerabilities, and timeline interfaces in benchmarks derived from the related literature, achieving 100\% recall and 81\% average precision.
These results showcase the potential of \tia to support diverse practical clients.
By open-sourcing all implementations supporting \tia's results, we seek to help advance static analysis for modern hardware designs.
\end{abstract}
